% Part: first-order-logic
% Chapter: natural-deduction
% Section: soundness-identity

\documentclass[../../../include/open-logic-section]{subfiles}

\begin{document}

\olfileid{fol}{ntd}{sid}

\olsection{صحت با \usetoken{S}{identity}}

\begin{prop}
استنتاج طبیعی با قواعد $\eq$ صحیح است.
\end{prop}

\begin{proof}
هر !!{formula} به صورت $\eq[t][t]$ معتبر است، زیرا برای هر
!!{structure}~$\Struct M$، $\Sat{M}{\eq[t][t]}$. (توجه کنید فرض
می‌کنیم ترم $t$ بسته است، یعنی هیچ متغیری ندارد؛ پس انتساب‌های متغیر
بی‌تأثیرند).

فرض کنید آخرین استنتاج در !!a{derivation}، \Elim{\eq} باشد؛ یعنی
!!{derivation} صورت زیر را دارد:
\begin{prooftree}
  \AxiomC{$\Gamma_1$}
  \RightLabel{$\delta_1$}
  \DeduceC{$\eq[t_1][t_2]$}
  \AxiomC{$\Gamma_2$}
  \RightLabel{$\delta_2$}
  \DeduceC{$!A(t_1)$}
  \RightLabel{\Elim{\eq}}
  \BinaryInfC{$!A(t_2)$}
\end{prooftree}
مقدمات $\eq[t_1][t_2]$ و $!A(t_1)$ به‌ترتیب ~!!{derive}d شده‌اند از
فرض‌های !!{undischarged}~$\Gamma_1$ و $\Gamma_2$. می‌خواهیم
نشان دهیم $!A(t_2)$ از $\Gamma_1 \cup \Gamma_2$ نتیجه می‌شود. یک
!!a{structure}~$\Struct{M}$ را با $\Sat{M}{\Gamma_1 \cup
  \Gamma_2}$ در نظر بگیرید. بنا بر فرض استقرا، $\Sat{M}{!A(t_1)}$ و
$\Sat{M}{\eq[t_1][t_2]}$. بنابراین
$\Value{t_1}{M} = \Value{t_2}{M}$. فرض کنید $s$ هر انتساب متغیری باشد و
$m = \Value{t_1}{M} = \Value{t_2}{M}$. بنا بر
\olref[fol][syn][ext]{prop:ext-formulas}، $\Sat{M}{!A(t_1)}[s]$ هرگاه و
تنها هرگاه $\Sat{M}{!A(x)}[\Subst{s}{m}{x}]$، هرگاه و تنها هرگاه
$\Sat{M}{!A(t_2)}[s]$. چون $\Sat{M}{!A(t_1)}$، داریم
$\Sat{M}{!A(t_2)}$.
\end{proof}

\end{document}
